\documentclass[11pt,a4paper,oneside]{book}
\usepackage[bindet]{90-faelles/latex/studiearkiv}
\hypersetup{pdftitle={Matematik – Bind 1: Algebra og lineær algebra}}
\setcounter{tocdepth}{2}
\begin{document}
\frontmatter
\hypersetup{pageanchor=false}
\begin{titlepage}
\centering
\vspace*{30mm}
{\sffamily\color{arkivblaa}\Large MATEMATIK TIL UNIVERSITETSFYSIK\par}
\vspace{16mm}
{\Huge\bfseries Bind 1\\[6mm]Algebra og lineær algebra\par}
\vspace{13mm}
{\Large Formler, begreber og matematisk sammenhæng\par}
\vfill
\begin{minipage}{0.8\textwidth}
Tal, funktioner, trigonometriske identiteter, vektorrum og lineære operatorer.

Dette bind genbruger de samme LaTeX-kapitler som det samlede matematikkompendium. Definitioner, metodevalg og formelopslag er samlet i tre kapitler. Der indgår ikke separate gennemregnede opgaveeksempler.
\end{minipage}\par
\vspace{14mm}
{\small Fysik-studiearkiv · Arbejdsudgave\par}
\end{titlepage}
\hypersetup{pageanchor=true}
\chapter*{Læsevejledning}
\addcontentsline{toc}{chapter}{Læsevejledning}
Dette bind er en emneopdelt udgave af det samlede matematikkompendium. Den generelle kilde og alle figurer vedligeholdes centralt under \texttt{00-matematik/chapters} og \texttt{00-matematik/figures}. For hele overblikket samt henvisninger mellem bindene anvendes den samlede udgave. Formelregisteret sidst i bogen giver direkte adgang til de centrale regler.
\section*{Lær, slå op, forbind}
Læs kapitlerne i rækkefølge for at bygge fra tal og funktioner til vektorer
og operatorer. Ved et konkret problem kan formelregisteret bruges som
indgang; læs altid betingelserne ved formlen, før den anvendes.

\hyperref[ch:algebra]{Algebra og funktioner} fastlægger domæner,
regneregler og grenvalg. \hyperref[ch:trigonometri]{Trigonometri} samler
vinkler, identiteter og faser. \hyperref[ch:b1-lineaer]{Lineær algebra}
samler ligningssystemer, geometri og normalmoder.

\section*{Notation}
\begin{tabular}{@{}lp{.68\textwidth}@{}}
\toprule
Symbol & Betydning\\\midrule
$\mathbb K$ & $\mathbb R$ eller $\mathbb C$, angivet ved behov\\
$\mathbb N,\mathbb N_0$ & Positive heltal, henholdsvis med nul\\
$A^{\mathsf T},A^\dagger$ & Transponeret, henholdsvis konjugeret transponeret\\
$\langle u,v\rangle$ & Indre produkt, lineært i andet argument\\
$\|v\|_2$ & Euklidisk norm i ortonormale koordinater\\
$\operatorname{rank},\ker,\operatorname{im}$ & Rang, kerne og billede\\
\bottomrule
\end{tabular}

Vinkler måles i radianer. Figurer bruger dimensionsløse koordinater,
medmindre en fysisk størrelse er angivet. Bindhenvisninger til senere
emner angiver deres placering i værket; de lover ikke, at de senere bind
er færdigskrevet. Ingen gennemregnede opgaveeksempler indgår.

\begingroup\small\tableofcontents\endgroup
\mainmatter
\chapter{Algebra og funktioner}
\label{ch:algebra}
Algebra er sproget, der forbinder en matematisk model med dens ukendte størrelser. I fysik er opgaven sjældent blot at isolere et symbol: Man skal også holde styr på, hvilke værdier det må antage, hvilke omskrivninger der er reversible, og om det endelige udtryk har de rigtige dimensioner. Kapitlet behandler derfor algebra som et \emph{opslagsværktøj} for senere modeller og ikke blot som en samling regneregler.

\section{Tal og mængder}
\label{sec:tal-maengder}
En mængde fastlægger, hvilke objekter en matematisk påstand handler om.
Dette er lige så vigtigt som selve regneoperationen: Et heltalligt antal
partikler, en reel position og en kompleks amplitude er forskellige slags
størrelser. En ligning beskriver derfor først et veldefineret problem, når
både de tilladte værdier og betingelserne for dens udtryk er angivet.

\subsection{Talområder og deres regnemuligheder}
\label{sec:talomraader}
Vi bruger følgende konventioner:
\begin{center}
\begin{tabular}{@{}p{0.16\textwidth}p{0.35\textwidth}p{0.39\textwidth}@{}}
\toprule
Talområde & Definition eller beskrivelse & Typisk fysisk rolle\\
\midrule
$\mathbb N$ & $\{1,2,3,\ldots\}$ & Positive tælleindekser\\[3pt]
$\mathbb N_0$ & $\{0,1,2,\ldots\}$ & Antal partikler eller besættelsestal\\[3pt]
$\mathbb Z$ & $\{\ldots,-1,0,1,\ldots\}$ & Heltallige modeindekser\\[3pt]
$\mathbb Q$ & $\{p/q:p\in\mathbb Z,\ q\in\mathbb Z\setminus\{0\}\}$ & Rationale forhold og endelige decimaltal\\[3pt]
$\mathbb R$ & Den fuldstændige ordnede talkrop & Kontinuerte koordinater og parametre\\[3pt]
$\mathbb C$ & $\{a+ib:a,b\in\mathbb R,\ i^2=-1\}$ & Amplituder og svingningsrepræsentationer\\
\bottomrule
\end{tabular}
\end{center}
Andre bøger kan medregne nul i $\mathbb N$; her bruges $\mathbb N_0$, når
nul skal indgå. Med de sædvanlige identifikationer har vi kæden
\[
\boxed{\mathbb N\subsetneq\mathbb N_0\subsetneq\mathbb Z
\subsetneq\mathbb Q\subsetneq\mathbb R\subsetneq\mathbb C.}
\]
Symbolet $\subsetneq$ betyder en ægte delmængde: alle elementer i den
venstre mængde ligger i den højre, men de to mængder er ikke ens.
Vi bruger $\subseteq$, når lighed er tilladt.

Et talområde er \emph{lukket} under en operation, hvis operationens
resultat igen ligger i området, når den er defineret. De naturlige tal er
lukkede under addition og multiplikation, men ikke under vilkårlig
subtraktion. De hele tal tillader også subtraktion, mens division kan kræve
rationale tal. Både $\mathbb Q$, $\mathbb R$ og $\mathbb C$ er talkroppe:
de fire regningsarter kan udføres inden for området, bortset fra division
med nul. Kvadratroden af et positivt rationalt tal behøver derimod ikke
være rational, så regningsarternes lukkethed dækker ikke alle funktioner.

De reelle tal har desuden en fuldstændighedsegenskab: Enhver ikke-tom
opadtil begrænset delmængde af $\mathbb R$ har en mindste øvre grænse i
$\mathbb R$. Den egenskab gør det muligt at begrunde grænseværdier og
konvergens i bind 2. At $\mathbb Q$ ligger tæt i $\mathbb R$ betyder, at
ethvert åbent interval med forskellige endepunkter indeholder et rationalt
tal; det betyder ikke, at alle reelle tal er rationale.

Komplekse tal udvider løsningsmulighederne: Ligningen $x^2+1=0$ har ingen
reel løsning, men løsningerne $\pm i$ i $\mathbb C$. De komplekse tal
har imidlertid ingen total
ordensrelation, der gør dem til en ordnet talkrop. Hvis en sådan orden
fandtes, ville kvadrater af ikke-nul tal være positive; så skulle både
$1=1^2$ og $-1=i^2$ være positive, hvilket er umuligt. En ulighed som
$z>0$ bruges derfor ikke for et generelt komplekst tal. Man kan i stedet
undersøge dets reelle del eller dets reelle modulus, afhængigt af spørgsmålet.

\subsection{Mængdenotation og logiske udsagn}
\label{sec:maengdenotation}
En mængde er bestemt af sine elementer; deres rækkefølge og eventuelle
gentagelser i en opremsning ændrer ikke mængden. Notationen $x\in A$
betyder, at $x$ tilhører $A$, mens $x\notin A$ betyder det modsatte.
Den tomme mængde $\varnothing$ har ingen elementer. Bemærk forskellen på
$\varnothing$ og $\{0\}$: den sidste mængde har ét element.

Mængdebyggernotationen
\[
A=\{x\in U:P(x)\}
\]
betyder alle de elementer i grundmængden $U$, der opfylder betingelsen
$P(x)$. Det lodrette skel læses »for hvilke«. En grundmængde kan være et
talområde, et rum af vektorer eller senere et rum af funktioner.
\begin{definition}[Delmængde og lighed]
$A\subseteq B$ betyder, at hvert element i $A$ også ligger i $B$.
To mængder er ens netop når $A\subseteq B$ og $B\subseteq A$.
\end{definition}
At vise mængdelighed kræver således to inklusioner. Det er en generel
bevismetode, der senere bruges ved sammenligning af kerner, billeder og
løsningsrum i lineær algebra.

Kvantorerne $\forall$ og $\exists$ betyder henholdsvis »for alle« og
»der findes«. Deres rækkefølge er afgørende: I
$\forall x\in A\;\exists y\in B:P(x,y)$ må valget af $y$ afhænge af $x$;
i $\exists y\in B\;\forall x\in A:P(x,y)$ skal det samme $y$ virke for
alle $x$. Denne forskel bliver central for punktvis og uniform konvergens
i bind 2. Negation af en kvantificeret påstand følger reglerne
\[
\boxed{\begin{aligned}
\neg\bigl(\forall x\in A:P(x)\bigr)
&\ \Longleftrightarrow\ \exists x\in A:\neg P(x),\\
\neg\bigl(\exists x\in A:P(x)\bigr)
&\ \Longleftrightarrow\ \forall x\in A:\neg P(x).
\end{aligned}}
\]
Et enkelt modeksempel afkræfter altså en universel påstand. Udsagnet
$P\Rightarrow Q$ siger, at $P$ er en tilstrækkelig betingelse for $Q$,
og at $Q$ er en nødvendig betingelse for $P$. En ækvivalens
$P\Leftrightarrow Q$ kræver begge implikationer. Denne skelnen beskytter
mod at forveksle en nødvendig kontrol med en fuldstændig karakterisering.

\subsection{Mængdeoperationer som kombination af betingelser}
\label{sec:maengdeoperationer}
Antag, at $A,B\subseteq U$, hvor $U$ er en fast grundmængde. De vigtigste
operationer er samlet her:
\begin{center}
\begin{tabular}{@{}p{0.22\textwidth}p{0.39\textwidth}p{0.29\textwidth}@{}}
\toprule
Operation & Elementbetingelse & Betydning\\
\midrule
$A\cap B$ & $x\in A$ og $x\in B$ & Begge krav gælder\\[3pt]
$A\cup B$ & $x\in A$ eller $x\in B$ & Mindst ét krav gælder\\[3pt]
$A\setminus B$ & $x\in A$ og $x\notin B$ & Udeluk elementer fra $B$\\[3pt]
$A^c=U\setminus A$ & $x\in U$ og $x\notin A$ & Komplement relativt til $U$\\
\bottomrule
\end{tabular}
\end{center}
»Eller« er her inklusivt: et element i begge mængder ligger også i deres
foreningsmængde. Komplementet afhænger af $U$, og grundmængden skal derfor
fremgå af sammenhængen.

De vigtigste omskrivninger er distributivitet og De Morgans love:
\[
\boxed{\begin{aligned}
A\cap(B\cup C)&=(A\cap B)\cup(A\cap C),\\
A\cup(B\cap C)&=(A\cup B)\cap(A\cup C),\\
(A\cup B)^c&=A^c\cap B^c,\\
(A\cap B)^c&=A^c\cup B^c.
\end{aligned}}
\]
Disse regler følger af betingelserne på hvert enkelt element. For den
første komplementregel gælder for eksempel
\[
x\in(A\cup B)^c
\ \Longleftrightarrow\ x\notin A\ \text{og}\ x\notin B
\ \Longleftrightarrow\ x\in A^c\cap B^c.
\]
Det er en udledning af en generel identitet, ikke en regel, der blot skal
huskes fra en tegning.

For en indekseret familie $(A_j)_{j\in J}$ betyder
$\bigcup_{j\in J}A_j$, at mindst én betingelse opfyldes, mens
$\bigcap_{j\in J}A_j$ kræver alle betingelser opfyldt. Når mange fysiske
bindinger skal gælde samtidig, er snitmængden derfor den naturlige
beskrivelse af de tilladte tilstande.

\subsection{Intervaller, endepunkter og ubegrænsede områder}
\label{sec:intervaller}
For reelle $a<b$ defineres
\[
\begin{aligned}
[a,b]&=\{x\in\mathbb R:a\leq x\leq b\},&
(a,b)&=\{x\in\mathbb R:a<x<b\},\\
[a,b)&=\{x\in\mathbb R:a\leq x<b\},&
(a,b]&=\{x\in\mathbb R:a<x\leq b\}.
\end{aligned}
\]
Runde parenteser udelader endepunktet, firkantede parenteser medtager det.
I dansk litteratur ses også $]a,b[$ for $(a,b)$. Her bruges den sidstnævnte
notation konsekvent. Et interval indeholder alle reelle tal mellem to af
sine punkter; en foreningsmængde af adskilte intervaller behøver ikke selv
være et interval.

Symbolerne $+\infty$ og $-\infty$ er ikke reelle tal og medtages aldrig
som intervalendepunkter:
\[
[a,\infty)=\{x\in\mathbb R:x\geq a\},\qquad
(-\infty,b)=\{x\in\mathbb R:x<b\}.
\]
En ubegrænset mængde beskriver vilkårligt store eller små værdier; den
indeholder ikke et særligt tal »uendelig«. Hvis endepunkterne falder sammen,
er $[a,a]=\{a\}$, mens $(a,a)=\varnothing$.

Figur~\ref{fig:algebra-intervaller} viser forskellen mellem medtagne og
udeladte punkter og mellem et interval og et område med et hul.

\begin{figure}[htbp]
\centering
\begin{tikzpicture}[x=0.92cm,y=0.85cm,>=stealth,
  every node/.style={font=\small}]
\foreach \y in {0,1.3,2.6} {
  \draw[->,gray] (-0.4,\y)--(7,\y) node[right] {$x$};
}
\draw[very thick,arkivblaa] (1,2.6)--(5.6,2.6);
\draw[thick,draw=arkivblaa,fill=white] (1,2.6) circle (2.7pt);
\fill[arkivblaa] (5.6,2.6) circle (2.7pt);
\node[below] at (1,2.5) {$a$};
\node[below] at (5.6,2.5) {$b$};
\node[anchor=east] at (-0.6,2.6) {$(a,b]$};
\draw[very thick,arkivblaa] (1,1.3)--(3.2,1.3);
\draw[very thick,arkivblaa] (3.2,1.3)--(5.6,1.3);
\foreach \x in {1,5.6} {
  \draw[thick,draw=arkivblaa,fill=white] (\x,1.3) circle (2.7pt);
}
\draw[thick,draw=arkivrod,fill=white] (3.2,1.3) circle (2.7pt);
\node[below] at (1,1.2) {$a$};
\node[below] at (3.2,1.2) {$c$};
\node[below] at (5.6,1.2) {$b$};
\node[anchor=east] at (-0.6,1.3) {$(a,b)\setminus\{c\}$};
\draw[very thick,arkivblaa] (1.4,0)--(5.2,0);
\foreach \x in {1.4,5.2} {
  \draw[thick,draw=arkivblaa,fill=white] (\x,0) circle (2.7pt);
}
\draw (3.3,-0.08)--(3.3,0.08);
\node[below] at (3.3,-0.1) {$x_0$};
\draw[<->,arkivrod] (3.3,0.4)--(5.2,0.4)
  node[midway,above] {$r$};
\node[below] at (1.4,-0.1) {$x_0-r$};
\node[below] at (5.2,-0.1) {$x_0+r$};
\node[anchor=east] at (-0.6,0) {$|x-x_0|<r$};
\end{tikzpicture}
\caption{Intervaller på en skematisk, dimensionsløs talakse. En fyldt prik
medtager et endepunkt, mens en åben prik udelader det. I midten antages
$a<c<b$: udeladelsen af $c$ deler området i to intervaller. Nederst vises
afstandsbetingelsen omkring $x_0$ med $r>0$.}
\label{fig:algebra-intervaller}
\end{figure}


Afstand på talaksen skrives $d(x,y)=|x-y|$. For en radius $r>0$ kan et
åbent interval derfor også beskrives som
\[
\boxed{\{x\in\mathbb R:|x-x_0|<r\}=(x_0-r,x_0+r).}
\]
Udledningen følger af $|u|<r\Leftrightarrow-r<u<r$ med $u=x-x_0$.
For $|x-x_0|\leq r$ medtages begge endepunkter. Denne notation beskriver
både tolerancer omkring en referenceværdi og de små omgivelser, der bruges
ved kontinuitet i bind 2. Absolutværdi og fejlgrænser behandles videre i
afsnittet om uligheder i dette kapitel.

\subsection{Produktmængder og flere variable}
\label{sec:produktmaengder}
Produktmængden $A\times B$ består af ordnede par:
\[
A\times B=\{(a,b):a\in A,\ b\in B\}.
\]
Rækkefølgen er en del af informationen. Generelt er $(a,b)$ ikke det
samme som $(b,a)$, selv om mængderne $\{a,b\}$ og $\{b,a\}$ er ens.
Notationen $\mathbb R^n$ betyder produktet af $n$ kopier af $\mathbb R$;
et element er en koordinatliste $(x_1,\ldots,x_n)$.

En produktmængde tillader uafhængige valg fra hver faktor. En binding
$F(x,y)=0$ udvælger derimod en delmængde af $\mathbb R^2$, hvor
koordinaterne ikke kan vælges frit. Dette er forudsætningen for senere at
beskrive løsningsmængder i lineær algebra, niveaukurver i bind 2 og
geometriske bindinger i bind 5. En koordinatliste alene fastlægger ikke et
fysisk vektorrum med enheder, norm og basis; den ekstra struktur skal
angives, når den bruges.

\subsection{Definitionsmængder og reversible omskrivninger}
\label{sec:definitionsmaengder}
Hvis et udtryk $F(x)$ skal være reelt, findes dets naturlige
\emph{definitionsmængde} ved at kræve, at alle dets dele er definerede.
Løsningsmængden for en ligning er derefter
\[
\boxed{S=\{x\in D:F(x)=0\},}
\]
hvor $D$ er den relevante definitionsmængde efter eventuelle yderligere
modelkrav. Betingelser skal kombineres med snitmængde, når alle dele af
udtrykket skal være definerede samtidigt:
\begin{center}
\begin{tabular}{@{}p{0.33\textwidth}p{0.57\textwidth}@{}}
\toprule
Reelt udtryk & Nødvendig definitionsbetingelse\\
\midrule
$P(x)/Q(x)$ & $Q(x)\ne0$\\[3pt]
$\sqrt{g(x)}$ & $g(x)\geq0$\\[3pt]
$1/\sqrt{g(x)}$ & $g(x)>0$\\[3pt]
$\ln g(x)$ & $g(x)>0$\\
\bottomrule
\end{tabular}
\end{center}
Desuden skal de indre udtryk være definerede. En bogstavelig formel kan
have et snævrere domæne end en forkortet repræsentation. Forkortning
\[
\frac{(x-a)P(x)}{x-a}=P(x)\qquad (x\ne a)
\]
beviser kun lighed på det oprindelige domæne. At tildele venstresiden en
værdi ved $x=a$ er en udvidelse af funktionen, ikke blot en omskrivning.

Division med et udtryk, der kan være nul, kan fjerne løsninger, mens
multiplikation med det kan tilføje løsninger. Kvadrering er heller ikke
generelt reversibel. For reelle $x,a$ gælder
\[
x=\sqrt a\quad\Longleftrightarrow\quad
\bigl(x^2=a\ \text{og}\ x\geq0\bigr)\qquad(a\geq0).
\]
Uden fortegnskravet på $x$ har man kun implikationen fra venstre mod
højre. Kandidater, der fremkommer ved en sådan omskrivning, skal derfor
kontrolleres i den oprindelige ligning. Det samme princip gælder, når
parameterafhængige nævnere eller inversgrene indgår.

\subsection{Metodevalg og fysisk gyldighed}
\label{sec:maengder-fysik}
Ved et algebraisk modelproblem er en nyttig rækkefølge at angive
grundmængden, opstille definitionsbetingelserne, skære dem med de fysiske
krav og først derefter løse ligningen. Ved flere tilladte alternativer
bruges en foreningsmængde. En parameter, hvis fortegn påvirker kravene,
skal behandles særskilt; et formelt resultat for positive parametre kan
ikke uden kontrol anvendes ved nul eller negative værdier.

For en endimensional konservativ bevægelse med konstant masse $m>0$ og
potentiel energi $V(q)$ giver energirelationen
\[
E=\frac{p^2}{2m}+V(q),\qquad
p^2=2m\bigl(E-V(q)\bigr).
\]
Et reelt momentum kræver $V(q)\leq E$. De energitilladte positioner er
således mængden
\[
A_E=\{q\in D_V:V(q)\leq E\}.
\]
Den kan bestå af flere adskilte intervaller; en partikel følger ikke
automatisk alle disse intervaller fra én begyndelsestilstand.
Ved et regulært vendepunkt med $V(q)=E$ er $p=0$. For $V(q)<E$ giver
$p=\pm\sqrt{2m(E-V(q))}$ to momentumgrene. Energiens nødvendige
positionsbetingelse vælger altså hverken bevægelsesretningen eller en
fuld tidsudvikling. Denne forskel mellem en tilladt mængde og en faktisk
løsning er vigtig i mekanik og i bind 3 om differentialligninger.

I statistisk fysik kan et besættelsestal ligge i $\mathbb N_0$, mens en
klassisk koordinat beskrives i $\mathbb R$. En kontinuert approksimation
til et diskret antal er en modelbeslutning og kræver sin egen begrundelse.
Mængdenotation holder disse valg synlige uden at foregribe den senere
fysiske teori.

\paragraph{Forbindelser i værket.}
Funktionsafsnittet nedenfor bruger domæne, codomæne og billede til at
præcisere, hvornår en invers eksisterer. Kapitlet om lineær algebra i dette
bind tilføjer regneoperationer og lineære bindinger på vektormængder.
Bind 2 bruger intervaller og kvantorer til grænseværdier, og bind 5
beskriver områder og koordinater i flere dimensioner. Disse henvisninger
angiver emneplacering; de er ikke en påstand om, at alle senere afsnit er
færdigskrevet.

\section{Potenser, rødder og brøker}
\label{sec:potenser-broeker}
Potenser samler gentagen multiplikation i ét symbol. Rødder vender denne
operation om, og brøker gør division til en del af regningen. De tre
værktøjer hænger tæt sammen, men deres regneregler har forskellige
forudsætninger. Særligt fortegn, nulværdier og ikke-heltallige eksponenter
kræver opmærksomhed, før en formel kan forkortes.

\subsection{Heltallige potenser og nulværdier}
For $n\in\mathbb N$ er $a^n$ produktet af $n$ faktorer lig med $a$.
For $a\ne0$ udvides definitionen ved
\[
a^0=1,\qquad a^{-n}=\frac{1}{a^n}.
\]
Med disse definitioner gælder for $m,n\in\mathbb Z$ og $a\ne0$
\[
\boxed{a^m a^n=a^{m+n},\qquad
\frac{a^m}{a^n}=a^{m-n},\qquad (a^m)^n=a^{mn}.}
\]
For positive eksponenter følger den første regel ved at tælle antallet af
faktorer. Kravet om at bevare samme regel ved $n=-m$ giver
$a^m a^{-m}=a^0=1$ og begrunder definitionen af negative potenser.
For $a=0$ er positive heltallige potenser definerede, men negative potenser
og $0^0$ bruges ikke her. En forkortning må derfor ikke skjule et
uoprindeligt defineret nuludtryk.

For heltallige $n$ gælder også
\[
(ab)^n=a^n b^n,\qquad
\left(\frac ab\right)^n=\frac{a^n}{b^n},
\]
hvor alle nævnere og negative potenser skal være definerede. En potens
fordeler sig over et produkt, men ikke generelt over en sum.
Fortegn og parenteser skal desuden holdes adskilt:
$(-a)^2=a^2$, mens $-a^2=-(a^2)$.

\subsection{Reelle rødder og rationale eksponenter}
For $n\in\mathbb N$ defineres den reelle $n$'te rod ved
\[
\begin{cases}
\sqrt[n]{a}=\text{det entydige }r\geq0\text{ med }r^n=a,
 & n\text{ lige},\ a\geq0,\\
\sqrt[n]{a}=\text{det entydige reelle }r\text{ med }r^n=a,
 & n\text{ ulige},\ a\in\mathbb R.
\end{cases}
\]
En lige rod er altså valgt ikke-negativ. Ligningen $x^2=a$ og selve
symbolet $\sqrt a$ har forskellige roller: for $a>0$ har ligningen to
rødder, men symbolet betegner kun den positive. Derfor gælder
\[
\boxed{\sqrt{x^2}=|x|\quad(x\in\mathbb R),\qquad
(\sqrt x)^2=x\quad(x\geq0).}
\]
Den første identitet følger, fordi $|x|$ er ikke-negativ og har kvadrat
$x^2$. At fjerne absolutværdien kræver den ekstra betingelse $x\geq0$.

For $a>0$ og heltal $p$, $q\in\mathbb N$ sætter man
\[
a^{p/q}=\bigl(\sqrt[q]{a}\bigr)^p.
\]
Værdien afhænger ikke af valget af brøkrepræsentation: Hvis $p/q=r/s$,
er de to positive kandidaters $qs$'te potenser begge $a^{ps}=a^{rq}$;
entydigheden af den positive rod giver samme værdi. Her må $p$ og $r$
også være negative, da $a\ne0$.
For $a=0$ kan positive rationale potenser defineres som nul.

For negative baser bruges en mere begrænset reel konvention:
Er $p/q$ forkortet med $q>0$ ulige, kan man sætte
$a^{p/q}=(\sqrt[q]{a})^p$. Ved lige nævner findes ingen reel værdi.
Man skal forkorte eksponenten først; ellers kan et valg af rod i en
mellemregning give et andet resultat. De generelle regler for sammensatte
potenser kan ikke overføres uden kontrol: Hvis $a<0$, er
$(a^2)^{1/2}=|a|\ne a^{2(1/2)}$.

\subsection{Vilkårlige reelle eksponenter og regneoversigt}
For en positiv base defineres vilkårlige reelle potenser ved
\[
\boxed{a^\alpha=\exp(\alpha\ln a)\qquad(a>0,\ \alpha\in\mathbb R).}
\]
Denne definition udvider de rationale potenser og giver de almindelige
potensregler uden grenproblemer. Reglerne følger af
$\exp(u+v)=\exp u\exp v$ og $\ln(ab)=\ln a+\ln b$ for positive
argumenter; eksponential- og logaritmefunktionerne behandles samlet
senere i dette kapitel.
\begin{center}
\begin{tabular}{@{}p{0.56\textwidth}p{0.34\textwidth}@{}}
\toprule
Regel & Tilstrækkelige reelle betingelser\\
\midrule
$a^\alpha a^\beta=a^{\alpha+\beta}$ & $a>0$\\[3pt]
$a^\alpha/a^\beta=a^{\alpha-\beta}$ & $a>0$\\[3pt]
$(a^\alpha)^\beta=a^{\alpha\beta}$ & $a>0$\\[3pt]
$(ab)^\alpha=a^\alpha b^\alpha$ & $a,b>0$\\[3pt]
$(a/b)^\alpha=a^\alpha/b^\alpha$ & $a,b>0$\\[3pt]
$\sqrt{ab}=\sqrt a\sqrt b$ & $a,b\geq0$\\[3pt]
$\sqrt{a/b}=\sqrt a/\sqrt b$ & $a\geq0,\ b>0$\\
\bottomrule
\end{tabular}
\end{center}
Betingelserne i tabellen er sikre fælles domæner; enkelte identiteter kan
være gyldige på større områder ved bestemte heltallige eksponenter.
Den reelle definition giver ikke vilkårlige irrationale potenser af
negative tal. Komplekse potenser kræver logaritmegrene og hører til bind 4.

\subsection{Brøkregning og fællesnævnere}
Brøken $a/b$ er for $b\ne0$ det tal, som multipliceret med $b$ giver $a$.
For reelle tal gælder
\[
\boxed{\begin{aligned}
\frac ab\pm\frac cd&=\frac{ad\pm bc}{bd}
&& (b,d\ne0),\\
\frac ab\,\frac cd&=\frac{ac}{bd}
&& (b,d\ne0),\\
\frac{a/b}{c/d}&=\frac{ad}{bc}
&& (b,c,d\ne0).
\end{aligned}}
\]
Additionsreglen udledes ved at skrive $a/b=ad/(bd)$ og
$c/d=bc/(bd)$. Man omskriver altså til samme enhed, før tællerne lægges
sammen. Divisionsreglen følger ved at gange med det reciprokke af den
ikke-nul divisor $c/d$.

Forkortning er division af hele tælleren og hele nævneren med en fælles
\emph{faktor}. Man kan ikke forkorte et enkelt led gennem en sum.
Hvis en fælles faktor afhænger af en variabel, skal dens nulpunkter
fortsat udelukkes fra det oprindelige domæne, jf.
afsnit~\ref{sec:definitionsmaengder}. Ved rationelle funktioner vælges en
fællesnævner efter faktorerne og deres nødvendige potenser; produktet af
alle nævnere virker, men kan være unødigt stort.

\subsection{Rationalisering og metodevalg}
Et rodudtryk i en forskel kan omskrives med den konjugerede faktor.
For $u,v\geq0$ og $u+v>0$ gælder
\[
\sqrt u-\sqrt v=\frac{u-v}{\sqrt u+\sqrt v}.
\]
Identiteten kommer af
$(\sqrt u-\sqrt v)(\sqrt u+\sqrt v)=u-v$.
Omvendt kan nævneren $\sqrt u-\sqrt v$ rationaliseres, når $u\ne v$:
\[
\frac{1}{\sqrt u-\sqrt v}
=\frac{\sqrt u+\sqrt v}{u-v}.
\]
Omskrivningerne er nyttige ved grænseværdier og ved forskelle mellem
næsten lige rødder. De udvider ikke automatisk domænet. Ved numerisk
beregning kan den første form reducere tabet af cifre i selve
rodsubtraktionen, men den samlede nøjagtighed afhænger også af, hvordan
$u-v$ beregnes; fejlanalyse hører til bind 4.

Vælg heltallige potensregler ved gentagne faktorer og reciprokke udtryk,
roddefinitionen ved spørgsmål om fortegn og reelle løsninger og positive
baser, når vilkårlige reelle eksponenter skal samles. Ved summer af brøker
findes fællesnævneren før forkortning. Ved rodligninger kontrolleres de
oprindelige fortegnskrav efter enhver kvadrering.

\subsection{Dimensionskontrol og fysiske skaleringslove}
Hvis $X$ er en fysisk størrelse med dimension $[X]$, har en heltallig
potens dimension $[X^n]=[X]^n$. Rødder og rationale potenser giver
tilsvarende brøkpotenser af enhederne, når størrelsen og den valgte
rodgren er meningsfulde. Addition kræver samme dimension i alle led,
og en algebraisk lighed kræver samme dimension på begge sider.

For en dimensionsbærende positiv størrelse er en skalering med en
vilkårlig eksponent klarest skrevet ved en positiv referenceskala $X_0$:
\[
\boxed{Y=Y_0\left(\frac{X}{X_0}\right)^\alpha.}
\]
Brøken $X/X_0$ er dimensionsløs, mens $Y_0$ har samme dimension som $Y$.
En potenslov $Y=CX^\alpha$ kan også bruges, men da afhænger dimensionen
af $C$ af $\alpha$. En variabel eksponent bør derfor ikke sættes på en
dimensionsbærende størrelse uden referenceskala.

En elastisk energi $U=kq^2/2$ med $k>0$ giver to mulige fortegn for
forskydningen, når energien og $k$ kendes. Kvadratroden bestemmer
størrelsen $|q|=\sqrt{2U/k}$, ikke bevægelsesretningen. Med
$[k]=\text{energi}/\text{længde}^2$ får roden dimension længde.
Dette forbinder rodens grenvalg med de tilladte tilstande fra
afsnit~\ref{sec:maengder-fysik}; den faktiske bevægelse kræver dynamik i
bind 3 og mekaniknoterne.

\section{Polynomier, faktorisering og binomialformlen}
\label{sec:polynomier}
Et polynomium bygger en funktion op af endeligt mange heltallige potenser.
Denne enkle form gør det muligt at samle led, finde nulpunkter og
omskrive en model uden at ændre dens indhold. Polynomier beskriver både
præcise algebraiske relationer og lokale tilnærmelser til mere komplicerede
funktioner. En faktorisering fremhæver rødderne, mens den udfoldede form
fremhæver koefficienterne; den nyttigste form afhænger af spørgsmålet.

\subsection{Koefficienter, grad og polynomiumsidentitet}
Et polynomium over $\mathbb K=\mathbb R$ eller $\mathbb C$ skrives
\[
P(x)=\sum_{k=0}^n a_k x^k=a_0+a_1x+\cdots+a_nx^n,
\qquad a_k\in\mathbb K.
\]
Hvis $a_n\ne0$, er graden $\deg P=n$, og $a_n$ er den ledende
koefficient. Et ikke-nul konstant polynomium har grad nul; nulpolynomiet
tilordnes ikke en endelig grad her. For ikke-nul polynomier gælder
\[
\deg(PQ)=\deg P+\deg Q.
\]
Dette følger, fordi produktets højeste led er
$a_n b_m x^{n+m}$ med $a_nb_m\ne0$. Derimod kan led annulleres ved
addition, så $\deg(P+Q)\leq\max(\deg P,\deg Q)$, hvis summen er ikke-nul.

To polynomier er identiske, når koefficienterne for hver potens er ens.
En identitet gælder for alle $x$ i talområdet; en ligning kan være sand
kun for bestemte værdier. Over $\mathbb R$ eller $\mathbb C$ giver lighed
som funktioner også lighed af koefficienter: et ikke-nul polynomium af
grad $n$ kan højst have $n$ forskellige rødder, som begrundes nedenfor.
Man kan derfor sammenligne koefficienter, når en omskrivning skal gælde
identisk. Et tjek i blot nogle få punkter beviser derimod ikke generelt
en identitet uden kontrol af graden.

\subsection{Standardfaktoriseringer og deres udledning}
Distributiviteten $a(b+c)=ab+ac$ giver følgende opslag:
\[
\boxed{\begin{aligned}
(a+b)^2&=a^2+2ab+b^2,\\
(a-b)^2&=a^2-2ab+b^2,\\
a^2-b^2&=(a-b)(a+b),\\
(a+b)^3&=a^3+3a^2b+3ab^2+b^3,\\
(a-b)^3&=a^3-3a^2b+3ab^2-b^3,\\
a^3-b^3&=(a-b)(a^2+ab+b^2),\\
a^3+b^3&=(a+b)(a^2-ab+b^2).
\end{aligned}}
\]
Kvadratsætningen fås ved at gange to faktorer $(a+b)$ ud; de to blandede
led er begge $ab$. Kubiksætningen følger ved at gange resultatet endnu
en gang med $(a+b)$. Det er netop de blandede led, der gør en potens af
en sum forskellig fra summen af potenser.

En generel forskel af potenser faktoriseres som
\[
a^n-b^n=(a-b)\sum_{k=0}^{n-1}a^{n-1-k}b^k
\qquad(n\in\mathbb N).
\]
Ved udmultiplikation annulleres alle mellemliggende led parvist.
Hvis $a\ne b$, kan man dividere med $a-b$; ved $a=b$ gælder den
uforkortede identitet stadig. Med $a=r$, $b=1$ giver samme idé
\[
\sum_{k=0}^{n-1}r^k=\frac{1-r^n}{1-r}\quad(r\ne1),
\qquad \sum_{k=0}^{n-1}1=n.
\]
Dette er en \emph{endelig} sum og kræver ingen konvergensbetingelse.
Uendelige geometriske rækker kræver en særskilt grænseovergang i bind 2.

\subsection{Binomialkoefficienter og binomialformlen}
Når produktet af $n$ faktorer $(a+b)$ udfoldes, opstår et led
$a^{n-k}b^k$ ved at vælge $b$ fra præcis $k$ af faktorerne. Antallet af
sådanne valg er binomialkoefficienten
\[
\binom nk=\frac{n!}{k!(n-k)!},\qquad
n\in\mathbb N_0,\quad k=0,\ldots,n,\quad 0!=1.
\]
Der er $n(n-1)\cdots(n-k+1)$ ordnede valg af $k$ forskellige positioner;
division med $k!$ fjerner rækkefølgen af de valgte positioner.
Dermed fås binomialformlen
\[
\boxed{(a+b)^n=\sum_{k=0}^n\binom nk a^{n-k}b^k
\qquad(n\in\mathbb N_0).}
\]
Koefficienterne er symmetriske, $\binom nk=\binom n{n-k}$, og opfylder
Pascals relation
\[
\binom nk=\binom{n-1}{k-1}+\binom{n-1}{k}\qquad(1\leq k\leq n-1).
\]
Et valg indeholder enten den sidste position eller gør det ikke; det
opdeler valgene i netop de to grupper. For $n=0$ forstås produktet som
et tomt produkt lig med 1; dette er ikke en generel definition af $0^0$
for vilkårlige reelle potenser.

Formlen gælder for almindelige reelle eller komplekse tal, hvor
multiplikationen er kommutativ. Den må ikke uden videre anvendes på
matricer: $(A+B)^2=A^2+AB+BA+B^2$ reduceres kun til den almindelige
kvadratsætning, hvis $AB=BA$. Potenser med ikke-heltallig eksponent giver
ikke denne endelige sum; binomialrækken behandles med sine
konvergensbetingelser i bind 2.

\subsection{Polynomiumsdivision, restsætning og rødder}
For polynomier $P,Q$ over $\mathbb R$ eller $\mathbb C$ med $Q\ne0$
findes entydige polynomier $S,R$, så
\[
\boxed{P(x)=Q(x)S(x)+R(x),\qquad R=0\ \text{eller}\ \deg R<\deg Q.}
\]
Metoden er at dividere det aktuelt højeste led i resten med det ledende
led i $Q$, føje kvotienten til $S$ og trække det tilsvarende multiplum af
$Q$ fra. Graden af den tilbageværende rest falder ved hvert trin, så
processen stopper. Hvis der fandtes to fremstillinger, ville
$Q(S_1-S_2)=R_2-R_1$ have grad mindst $\deg Q$ på venstresiden og mindre
på højresiden, medmindre begge sider er nul. Det begrunder entydigheden.

Ved division med $x-a$ er resten konstant. Indsættelse af $x=a$ giver
restsætningen og faktorsætningen:
\[
\boxed{P(x)=(x-a)S(x)+P(a),\qquad
P(a)=0\ \Longleftrightarrow\ (x-a)\text{ er faktor i }P.}
\]
Et ikke-nul polynomium af grad $n$ har derfor højst $n$ forskellige
rødder: Hver forskellig rod fjerner en lineær faktor, og graden falder
én ad gangen. En rod $a$ har multiplicitet $m$, hvis
$P(x)=(x-a)^m H(x)$ med $H(a)\ne0$. Multipliciteten er altså antallet af
gange faktoren forekommer, ikke antallet af forskellige rodværdier.

Algebraens fundamentalsætning siger, at et polynomium af grad $n\geq1$
over $\mathbb C$ kan faktoriseres i $n$ lineære faktorer talt med
multiplicitet. Sætningen anføres her uden bevis. Har polynomiet reelle
koefficienter, gælder $P(\bar z)=\overline{P(z)}$, så ikke-reelle rødder
kommer i konjugerede par. Et par $a\pm ib$ med $b\ne0$ giver faktoren
\[
\bigl(x-(a+ib)\bigr)\bigl(x-(a-ib)\bigr)
=(x-a)^2+b^2,
\]
som er irreducibel over $\mathbb R$. Et reelt polynomium kan derfor
faktoriseres i reelle lineære og irreducible kvadratiske faktorer.
Det er grundlaget for partialbrøksopdelingen senere i kapitlet.

\subsection{Metodevalg, fortegn og fysiske forbindelser}
Faktorisér, når nulpunkter, fælles faktorer eller fortegn er det centrale;
udfold og sammenlign koefficienter, når en identitet eller en
parameterbestemmelse skal kontrolleres. Brug polynomiumsdivision, når en
brøks tæller har mindst samme grad som nævneren, eller når et kendt
nulpunkt skal skilles ud. En eksakt faktoriseringsmulighed er nyttig,
men der findes ikke en generel rodformel med radikaler for alle
polynomier af grad fem eller højere; numeriske rødder behandles i bind 4.

Omkring en reel rod $a$ med multiplicitet $m$ er fortegnet af
$P(x)=(x-a)^m H(x)$ bestemt af $x-a$ og det lokalt uændrede fortegn af
$H$. Derfor skifter $P$ fortegn gennem roden ved ulige $m$, men ikke ved
lige $m$. Denne sammenhæng mellem faktorisering og grafen er vist i
figur~\ref{fig:algebra-multiplicitet} og bruges ved uligheder og ved
undersøgelse af, hvor en fysisk størrelse kan være positiv.

\begin{figure}[htbp]
\centering
\begin{tikzpicture}
\begin{axis}[width=0.79\textwidth,height=6cm,
 axis lines=middle,xmin=-1.15,xmax=1.15,ymin=-1.15,ymax=1.3,
 xlabel={$u=x-a$},ylabel={$P$},
 xtick={-1,0,1},ytick={-1,1},
 samples=120,domain=-1.05:1.05,
 legend style={at={(0.03,0.97)},anchor=north west,
 draw=none,fill=white,font=\small},
 label style={font=\small}]
\addplot[arkivblaa,very thick] {x^2};
\addlegendentry{$P=u^2$: lige multiplicitet}
\addplot[arkivrod,very thick,dashed] {x^3};
\addlegendentry{$P=u^3$: ulige multiplicitet}
\end{axis}
\end{tikzpicture}
\caption{Dimensionsløse grundformer ved en rod $x=a$.
En rod med lige multiplicitet berører nul uden fortegnsskift;
en rod med ulige multiplicitet krydser nul. For et generelt polynomium
$(x-a)^mH(x)$ gælder fortegnsargumentet lokalt, hvor $H$ er forskellig
fra nul; $H$ kan ændre både skala og kurvens orientering.}
\label{fig:algebra-multiplicitet}
\end{figure}


For en dimensionsbærende koordinat $q$ skal koefficienterne i en energi
$V(q)=\sum_k c_k q^k$ opfylde $[c_k]=[\text{energi}]/[q]^k$.
Ved symmetrien $V(-q)=V(q)$ for alle $q$ må de ulige koefficienter være
nul; dette følger ved at sammenligne koefficienter. Et lige potentiale
kan derfor organiseres med lige potenser uden at gætte dets fysiske
parametre. En lokal polynomiumsapproksimation er derimod ikke en
identitet for hele potentialet; dens restled og gyldighedsinterval
behandles i bind 2 om Taylorudvikling.

I lineær algebra er $\det(\lambda I-A)$ et polynomium i $\lambda$.
Faktorisering og multiplicitet vender derfor tilbage ved egenværdier,
men en gentagen polynomiumsrod garanterer ikke flere uafhængige
egenvektorer. Den ekstra geometriske betingelse behandles i kapitlet om
lineær algebra i dette bind.

\section{Ligninger og uligheder}
\subsection{Andengradsligninger og kvadratkomplettering}
Den generelle andengradsligning $ax^2+bx+c=0$ med $a\ne0$ har rødder
\[
\boxed{x_\pm=\frac{-b\pm\sqrt{b^2-4ac}}{2a}}.
\]
Diskriminanten $\Delta=b^2-4ac$ afgør antallet af reelle rødder. Fuldstændiggørelse af kvadratet giver
\[
ax^2+bx+c=a\left(x+\frac{b}{2a}\right)^2+
c-\frac{b^2}{4a}.
\]
Denne form er ofte mere informativ end rodformlen, når et minimum af en energifunktion eller en kvadratisk approksimation skal bestemmes.

\subsection{Brøker og partialbrøker}
En rational funktion $P(x)/Q(x)$ kan ofte dekomponeres efter polerne i nævneren. For $a\ne b$ er eksempelvis
\[
\frac{1}{(x-a)(x-b)}
=\frac{1}{a-b}\left(\frac{1}{x-a}-\frac{1}{x-b}\right).
\]
Gentagne lineære faktorer $(x-a)^m$ kræver led med alle nævnerpotenser $(x-a)^{-k}$ for $1\leq k\leq m$. Irreducible reelle andengradsfaktorer kræver en lineær tæller. Hvis tællerens grad ikke er mindre end nævnerens, udføres først polynomiumsdivision. Metoden bruges især før integration eller Laplace-inversion.

\subsection{Uligheder og absolutværdi}
Absolutværdi er afstanden fra nul på den reelle talakse:
\[
|x|=\begin{cases}x,&x\geq0,\\-x,&x<0.\end{cases}
\qquad |x+y|\leq |x|+|y|.
\]
Hvis en ulighed multipliceres med et negativt tal, skifter ulighedstegnet retning. Ved multiplikation med en ukendt parameter skal man derfor skelne mellem dens fortegnstilfælde.

For $a>0$ gælder den nyttige ækvivalens
\[
|x-x_0|<a\quad\Longleftrightarrow\quad x_0-a<x<x_0+a.
\]
I perturbationsteori giver en antagelse som $|\varepsilon|\ll1$ ikke alene en streng matematisk fejlgrænse. En kontrolleret tilnærmelse skal kobles til et restled eller en eksplicit størrelsesorden.

\section{Funktioner, inverse funktioner og sammensætning}
\subsection{Domæne, billede og sammensætning}
En funktion $f:D\to Y$ tildeler hvert $x\in D$ præcis én værdi $f(x)\in Y$. Definitionsmængden $D$, codomænet $Y$ og billedmængden $f(D)$ er forskellige begreber. En fysisk formel kan være algebraisk meningsfuld for et større interval end det, hvor modellen faktisk gælder.

Sammensætningen $f\circ g$ defineres ved $(f\circ g)(x)=f(g(x))$, når $x$ ligger i $g$'s definitionsmængde, og $g(x)$ desuden ligger i $f$'s definitionsmængde. Domænet for den sammensatte funktion er derfor
\[
D_{f\circ g}=\{x\in D_g:g(x)\in D_f\}.
\]
Et vigtigt modelspørgsmål er, om en funktion er injektiv (forskellige input giver forskellige output), surjektiv (hele codomænet rammes) eller bijektiv (begge dele). Kun en bijektiv funktion har en entydig global invers som afbildning mellem de angivne mængder.

\subsection{Inverse funktioner og implicitte relationer}
Hvis $f$ er strengt monoton på et interval, er den injektiv dér og har en invers $f^{-1}:f(D)\to D$. Hvis $f$ er differentiabel, og $f'(x)\ne0$, gælder lokalt
\[
\boxed{(f^{-1})'(f(x))=\frac{1}{f'(x)}}.
\]
Antagelsen $f'(x)\ne0$ er afgørende. For eksempel har $f(x)=x^3$ en global invers, men formlen giver ingen endelig afledt for inversen ved $f(0)=0$.

En implicit relation $F(x,y)=0$ beskriver ikke nødvendigvis en global funktion $y(x)$. Hvis $F$ er kontinuert differentiabel og $F_y(x_0,y_0)\ne0$, garanterer den implicitte funktionssætning en lokal gren med
\[
\frac{\dd y}{\dd x}=-\frac{F_x}{F_y}.
\]
Det er et typisk værktøj, når en bane eller en bevaret energi definerer en relation mellem størrelser, som er besværlig at løse eksplicit.

\subsection{Skalering, dimensionsløse variable og følsomhed}
En størrelse $x$ med dimension kan gøres dimensionsløs ved $u=x/x_0$, hvor $x_0$ er en referenceværdi med samme dimension. Hvis en model har formen $y=A\,f(x/x_0)$, beskriver $f$ en universel form, mens $A$ og $x_0$ sætter de fysiske skalaer.

For en positiv, differentiabel størrelse $y=f(x)$ med $x>0$ er den logaritmiske følsomhed
\[
\frac{\dd\ln y}{\dd\ln x}=\frac{x}{y}\frac{\dd y}{\dd x}.
\]
Den måler en relativ ændring i output pr. relativ ændring i input. For en potenslov $y=Cx^\alpha$ er følsomheden netop $\alpha$. Det hjælper med at afgøre, hvilke inputusikkerheder der er mest betydningsfulde, men usikkerhedsforplantning kræver yderligere antagelser om fejl og korrelationer.

\section{Eksponential- og logaritmefunktioner}
\subsection{Eksponentiel vækst og henfald}
For $a>0$ defineres $a^x=e^{x\ln a}$. Eksponentialfunktionen er positiv og opfylder
\[
e^{x+y}=e^x e^y,\qquad
\frac{\dd}{\dd x}e^{kx}=k e^{kx}.
\]
Differentialligningen $\dot y=ky$ har løsningen $y(t)=y(0)e^{kt}$ for konstant $k$. Her har $k$ dimension invers tid. Ved $k<0$ beskriver modellen eksponentielt henfald, og karakteristisk tid er $\tau=1/|k|$. Ved $k>0$ beskriver den eksponentiel vækst; fordoblingstiden er $\ln2/k$.

Logaritmen er invers til eksponentialfunktionen på de positive reelle tal. For $x,y>0$ og reel $\alpha$ gælder
\[
\ln(xy)=\ln x+\ln y,\qquad
\ln\frac{x}{y}=\ln x-\ln y,\qquad
\ln(x^\alpha)=\alpha\ln x.
\]
Man må ikke skrive $\ln(x+y)=\ln x+\ln y$. Logaritmen omdanner produkter til summer, ikke addition til en simpel operation.

% Original principskitse til algebra: lineær, eksponentiel og logaritmisk skalering.
\begin{figure}[htbp]
\centering
\begin{tikzpicture}
\begin{axis}[
 width=0.86\textwidth,height=6.3cm,
 axis lines=middle,xmin=-1.3,xmax=3.3,ymin=-1.4,ymax=4.6,
 samples=170,domain=-1.2:3.2,
 xlabel={$x$},ylabel={$y$},
 grid=major,grid style={gray!15},
 legend style={at={(0.98,0.03)},anchor=south east,fill=white,draw=none,font=\small}
]
\addplot[very thick,color=arkivblaa] {x};
\addlegendentry{$y=x$}
\addplot[very thick,color=arkivrod] {exp(x)};
\addlegendentry{$y=e^x$}
\addplot[very thick,color=black,domain=0.06:3.2] {ln(x)};
\addlegendentry{$y=\ln x$}
\addplot[only marks,mark=*,mark size=1.5pt,color=arkivrod] coordinates {(0,1)};
\addplot[only marks,mark=*,mark size=1.5pt,color=black] coordinates {(1,0)};
\end{axis}
\end{tikzpicture}
\caption{De tre forskellige former for afhængighed, vist med dimensionsløse variable: lineær vækst, eksponentialfunktion og logaritme. Eksponentialfunktionen og logaritmen er hinandens inverse efter spejling i linjen $y=x$. Grafen viser kun et endeligt udsnit af de respektive definitionsområder.}
\label{fig:algebra-eksplog}
\end{figure}


\subsection{Logaritmer, grene og fysiske enheder}
I en fysisk formel skal et logaritme- eller eksponentialargument være dimensionsløst. Udtrykket $\ln x$ er derfor ikke meningsfuldt, hvis $x$ f.eks. betegner en længde uden en angivet enhedsskala. Brug i stedet $\ln(x/x_0)$ for en positiv referencelængde $x_0$.

For $x>0$ gælder den nyttige relation
\[
\frac{\dd}{\dd x}\ln x=\frac{1}{x},
\]
og for $x\ne0$ er den reelle afledte af $\ln|x|$ også $1/x$. I kompleks analyse bliver logaritmen flerbladet, så en gren skal vælges særskilt; de almindelige reelle logaritmeregler må ikke overføres ukritisk til vilkårlige komplekse argumenter.

\section{Metodevalg, kontrol og anvendelser i fysik}
Når en fremtidig mekanik- eller elektromagnetismeopgave reduceres til algebra, kan følgende rækkefølge genbruges:
\begin{enumerate}
\item Bestem de ukendte størrelser, deres dimensioner og definitionsmængder.
\item Identificér mulige symmetrier, bevarede størrelser og dimensionsløse parametre, før du isolerer variable.
\item Vælg en algebraisk omskrivning, der bevarer ligningens løsninger; noter særskilt, når dette ikke er garanteret.
\item Brug en relevant repræsentation: faktorisering til rødder, fuldstændigt kvadrat til ekstrema, partialbrøker til integration eller logaritmer til produkter og skalering.
\item Indsæt resultatet i den oprindelige ligning, og kontrollér grænsetilfælde og enheder.
\end{enumerate}
For et central-kraftproblem bruges f.eks. løsning af en andengradsligning til radialvendepunkter, og fuldstændiggørelse af kvadratet kan hjælpe med at undersøge et potentiales form. I lineær mekanik kræver et stort koblet ligningssystem ofte matricemetoder snarere end manuel elimination; disse behandles i kapitlet om lineær algebra.

\paragraph{Opslag: hyppige fejl.}
Undgå division med mulige nulled, forkert fortegn ved uligheder, skjulte inversgrene, eksponentialfunktioner med dimensionsbærende argumenter og antagelser om, at en lokal linearisering gælder globalt. Hvis modellen introducerer en parameter, bør særlige værdier som nul og grænsetilfælde behandles eksplicit.

\chapter{Trigonometri}
\planlagt{Trigonometriske definitioner, identiteter, inverser og hyperbolske funktioner udvikles samlet.}
\section{Trigonometriske identiteter og hyperbolske funktioner}
\leavevmode\par
% Redaktionelt: tilføj definitioner, gyldighed, relationer og relevante figurer.

\chapter{Lineær algebra}
Lineær algebra beskriver vektorrum og lineære afbildninger mellem dem. Dens styrke er, at mange problemer inden for mekanik, elektromagnetisme, differentialligninger og kvantemekanik kan samles i den samme struktur: vektorer, operatorer og deres invariant definerede egenskaber. En matrix er en \emph{repræsentation} af en lineær afbildning i valgte baser; afbildningen selv afhænger ikke af basisvalget.

\section{Vektorrum, basis og lineære afbildninger}
Et vektorrum $V$ over et legeme $\mathbb K$ (typisk $\mathbb R$ eller $\mathbb C$) er lukket under addition og skalarmultiplikation og opfylder vektorrumsaksiomerne. Et sæt $(v_1,\ldots,v_n)$ er en basis, hvis enhver $v\in V$ har en \emph{entydig} fremstilling
\[
v=\sum_{i=1}^{n}v^i v_i.
\]
Koefficienterne $v^i$ er koordinater i den valgte basis. De er ikke selve vektoren. I et endeligdimensionalt rum har alle baser samme antal elementer, nemlig $\dim V$.

En afbildning $T:V\to W$ er lineær, når
\[
T(av+bw)=aT(v)+bT(w)
\qquad(a,b\in\mathbb K).
\]
I baserne $(v_j)$ og $(w_i)$ defineres matricen $A=(A^i{}_j)$ ved
\[
T(v_j)=\sum_i A^i{}_j w_i,\qquad [T(v)]_w=A[v]_v.
\]
Ved et basisskift $[v]_{\mathrm{gammel}}=S[v]_{\mathrm{ny}}$ for en operator på samme rum gælder
\[
A_{\mathrm{ny}}=S^{-1}A_{\mathrm{gammel}}S.
\]
Dette er en \emph{similaritetstransformation}. Den bevarer karakteristisk polynomium, egenværdier, determinant og spor. Den må ikke forveksles med omregning af en bilineær form, hvor transformationsloven indeholder en transponeret matrix.

\subsection{Kerne, billede og dimension}
Kernen og billedet defineres ved
\[
\ker T=\{v\in V:T(v)=0\},\qquad
\operatorname{im}T=\{T(v):v\in V\}.
\]
For endeligdimensionalt $V$ giver dimensionssætningen
\[
\boxed{\dim V=\dim(\ker T)+\operatorname{rank}T.}
\]
En lineær afbildning er injektiv præcis når dens kerne kun består af nulvektoren. For en kvadratisk matrix er injektivitet, surjektivitet og invertibilitet ækvivalente.

\section{Determinanter, rang og lineære ligninger}
Et lineært ligningssystem $Ax=b$ er konsistent, hvis og kun hvis $b$ ligger i søjlerummet for $A$, eller ækvivalent
\[
\operatorname{rank}(A)=\operatorname{rank}([A\mid b]).
\]
Hvis løsninger findes, er deres samlede form $x=x_{\mathrm p}+z$, hvor $Az=0$. Antallet af frie parametre er $\dim\ker A$. Gauss-elimination ændrer ikke løsningsmængden, når man udfører elementære rækkeoperationer på den udvidede matrix.

Determinanten findes for kvadratiske matricer. Dens absolutværdi er volumenfaktoren for den tilsvarende lineære transformation, mens fortegnet i reelle rum angiver orienteringsskift:
\[
\det(AB)=\det A\det B,\qquad
\det(A^{\mathsf T})=\det A,\qquad
\det(A^{-1})=\frac{1}{\det A}\quad(\det A\neq0).
\]
En kvadratisk matrix er invertibel præcis når dens determinant er forskellig fra nul. Determinanten kan dog være numerisk uhensigtsmæssig som praktisk test for næsten-singularitet: konditionstal og singularværdier giver ofte en bedre vurdering.

\subsection{Spor og karakteristisk polynomium}
Sporet er summen af diagonalelementerne og er basisuafhængigt:
\[
\operatorname{tr}(A+B)=\operatorname{tr}A+\operatorname{tr}B,\qquad
\operatorname{tr}(AB)=\operatorname{tr}(BA).
\]
For en $n\times n$-matrix er det karakteristiske polynomium
\[
p_A(\lambda)=\det(\lambda I-A).
\]
Over $\mathbb C$ er summen af egenværdierne med algebraisk multiplicitet $\operatorname{tr}A$, og produktet er $\det A$.


\section{Fra diagonalelementer til spor og tensorkontraktion}
\label{sec:spor-kontraktion}
En vigtig idé, som vender tilbage i tensorregning og differentialgeometri, er at danne en skalar ved at summere over et par sammenhørende indekser. I lineær algebra møder vi den først som \emph{sporet} af en lineær operator.

For en kvadratisk matrix $A=(A^i{}_j)$ er
\[
\boxed{\operatorname{tr}A=\sum_{i=1}^{n}A^i{}_i.}
\]
Summen løber over \emph{diagonalelementerne}, men sporet er ikke det samme som at finde eller konstruere en diagonal matrix. Som en første illustration er
\[
A=\begin{pmatrix}2&5\\-1&3\end{pmatrix},
\qquad
\operatorname{tr}A=2+3=5.
\]
Derimod betyder \emph{diagonalisering} at finde en basis, hvor matrixrepræsentationen kan skrives $D=P^{-1}AP$ med alle elementer uden for diagonalen lig nul. Dette kræver tilstrækkeligt mange lineært uafhængige egenvektorer.

Hvorfor er sporet uafhængigt af basis? Under et basisskift bliver
$A'=S^{-1}AS$. Da $\operatorname{tr}(BC)=\operatorname{tr}(CB)$ for kvadratiske matricer med kompatible dimensioner, får vi
\[
\operatorname{tr}(S^{-1}AS)
=\operatorname{tr}(ASS^{-1})
=\operatorname{tr}A.
\]
Den specifikke diagonal ændrer sig altså ved basisskift, men \emph{summen} af dens elementer gør ikke.

I tensorregning generaliseres dette til \emph{kontraktion} af et øvre og et nedre indeks, f.eks.
\[
T^i{}_i=\sum_i T^i{}_i.
\]
I dette eksempel er kontraktionen netop sporet. For en tensor med flere indekser kan én kontraktion efterlade en tensor af lavere orden. Man kan ikke generelt kontrahere to vilkårlige øvre indekser uden yderligere struktur: en metrik kan bruges til at hæve eller sænke indeks, så parringen giver geometrisk mening.

\paragraph{Forbindelse fremad.}
I bind 5, kapitlet om tensorregning, bruges samme idé til f.eks.
$R=g^{ij}R_{ij}$: Ricciskalaren dannes ved at kontrahere Riccitensoren med den inverse metrik. Sporet af en matrix er derfor et første, konkret tilfælde af en langt bredere geometrisk operation.

\paragraph{Hvornår bruger man hvad?}
Vil man beregne summen af egenværdierne med algebraisk multiplicitet, bruger man sporet. Vil man finde de retninger, hvor en lineær operator virker ved simpel skalering, undersøger man egenvektorer og diagonalisering. Vil man reducere antallet af frie indekser i en tensorlig relation på en koordinatuafhængig måde, bruger man kontraktion.

\section{Egenværdier og diagonalisering}
En ikke-nul vektor $v$ kaldes egenvektor for $A$, hvis $Av=\lambda v$. Egenværdierne opfylder
\[
\boxed{\det(A-\lambda I)=0.}
\]
Egenvektorer med indbyrdes forskellige egenværdier er lineært uafhængige. Derfor kan en $n\times n$-matrix diagonaliseres, hvis den har $n$ lineært uafhængige egenvektorer:
\[
A=PDP^{-1},\qquad D=\operatorname{diag}(\lambda_1,\ldots,\lambda_n).
\]
Det er ikke tilstrækkeligt blot at have $n$ egenværdier \emph{talt med algebraisk multiplicitet}: en matrix kan mangle egenvektorer. For en sådan matrix kan Jordan-normalform være relevant.

Når diagonalisering findes, kan man beregne funktioner af matricen ved at anvende funktionen på diagonalelementerne:
\[
f(A)=P\,\operatorname{diag}\bigl(f(\lambda_1),\ldots,f(\lambda_n)\bigr)P^{-1}.
\]
For den lineære differentialligning $\dot x=Ax$ giver dette særligt
\[
x(t)=e^{At}x(0),\qquad
e^{At}=Pe^{Dt}P^{-1}.
\]
Matrixeksponentialen er også defineret, når $A$ ikke kan diagonaliseres, ved sin absolut konvergente potensrække.

\section{Indre produkter, spektralsætning og SVD}
Et indre produkt $\langle u,v\rangle$ giver et mål for længde og ortogonalitet. På $\mathbb R^n$ bruges typisk $u^{\mathsf T}v$ og på $\mathbb C^n$ det hermiteske produkt $u^\dagger v$ (konjugeret lineært i første argument). Normen er
\[
\|v\|=\sqrt{\langle v,v\rangle},\qquad
|\langle u,v\rangle|\leq \|u\|\|v\|.
\]
En ortonormal basis opfylder $\langle e_i,e_j\rangle=\delta_{ij}$. I en sådan basis er den adjungerede operators matrix $A^\dagger$.

\subsection{Selvadjungerede og normale operatorer}
En kompleks matrix er hermitesk, hvis $A^\dagger=A$, og en reel matrix er symmetrisk, hvis $A^{\mathsf T}=A$. Den endeligdimensionale spektralsætning siger, at en normal kompleks matrix ($AA^\dagger=A^\dagger A$) kan diagonaliseres ved en unitær matrix:
\[
A=UDU^\dagger,\qquad U^\dagger U=I.
\]
For en hermitesk matrix er egenværdierne reelle. Egenvektorer svarende til forskellige egenværdier er ortogonale; dette er centralt for blandt andet normale svingninger og kvantemekaniske observabler.

\subsection{Singularværdidekomposition}
Enhver matrix $A\in\mathbb C^{m\times n}$ har en singularværdidekomposition
\[
\boxed{A=U\Sigma V^\dagger,}
\]
hvor $U,V$ er unitære, og $\Sigma$ er rektangulær diagonal med ikke-negative singularværdier $\sigma_i$. De ikke-nul singularværdier er kvadratrødderne af de positive egenværdier for $A^\dagger A$. Antallet af positive singularværdier er matricens rang.

SVD findes også, når matricen ikke er kvadratisk eller diagonaliserbar. Den er derfor et generelt værktøj til projektion, rangvurdering og mindste-kvadraters problemer. For en matrix med fuld søjlerang er mindste-kvadraters løsningen til $\min_x\|Ax-b\|_2$ entydig; uden fuld rang vælges typisk den løsning, der har mindst norm, ved hjælp af pseudoinversen $A^+$.

\subsection{Ortogonale projektioner}
For en ortonormal søjlematrix $Q$ med $Q^\dagger Q=I$ er
\[
P=QQ^\dagger,\qquad P^2=P,\qquad P^\dagger=P.
\]
$P$ projicerer ortogonalt på søjlerummet for $Q$. Det gør relationen mellem lineær algebra, geometri og approksimation tydelig: en bedste approksimation i et underrum fås ved at fjerne den ortogonale restkomponent.

\paragraph{Anvendelser i fysik.}
Normalmoders frekvenser findes typisk fra generaliserede egenværdiproblemer; lineæriserede bevægelsesligningers stabilitet bestemmes blandt andet af operatorens spektrum; ortonormale baser og selvadjungerede operatorer er centrale i kvantemekanik. For ikke-ortogonale koordinater skal metrikken indgå eksplicit i indre produkter og adjungering.

\section{Praktisk metodevalg: lineære systemer og egenværdier}
\label{sec:la-metode}
Inden man regner, bør man skelne mellem tre forskellige spørgsmål. Et ligningssystem $Ax=b$ handler om en \emph{løsning}; egenværdiproblemet $Av=\lambda v$ handler om operatorens \emph{naturlige retninger}; og et mindste-kvadraters problem handler om den \emph{bedst mulige approksimation}, når ligningerne ikke kan opfyldes præcist.

\subsection{Sådan vælges metoden}
\begin{center}
\begin{tabular}{@{}p{0.32\textwidth}p{0.57\textwidth}@{}}
\toprule
Spørgsmål & Typisk fremgangsmåde\\
\midrule
Er $Ax=b$ løsbart? & Sammenlign $\operatorname{rank}A$ og $\operatorname{rank}[A\mid b]$.\\
Er løsningen entydig? & Kræv en konsistent ligning og $\ker A=\{0\}$.\\
Er en kvadratisk matrix invertibel? & Brug rang eller determinant; numerisk også singularværdier.\\
Hvilke naturlige moder har systemet? & Stil egenværdiproblemet op, kontroller massen og fortegnene.\\
Kan operatoren diagonaliseres? & Kontroller antal lineært uafhængige egenvektorer; ikke blot antal rødder.\\
Hvilken approximation er nærmest $b$? & Benyt ortogonal projektion eller SVD.\\
\bottomrule
\end{tabular}
\end{center}
I numeriske beregninger bør man normalt løse $Ax=b$ med en faktorisering eller eliminationsmetode i stedet for først at beregne $A^{-1}$. At inversen findes matematisk, siger ikke alene at problemet er numerisk stabilt.

\subsection{Matrix og koordinatvalg}
Lad $v$ være en geometrisk vektor. Den kan have forskellige koordinatsøjler $[v]_e$ og $[v]_f$ i to baser. Skriver vi $[v]_e=S[v]_f$, følger for en operator $T:V\to V$
\[
[T]_f=S^{-1}[T]_eS.
\]
Denne formel forklarer, hvorfor eigenværdier og spor er basisuafhængige, selv om alle de enkelte matrixelementer kan ændre sig. Hvis man i stedet skifter baser på både domæne og codomæne i en afbildning $T:V\to W$, bliver den generelle regel $[T]_{\mathrm{ny}}=R^{-1}[T]_{\mathrm{gammel}}S$, hvor $S$ og $R$ beskriver skiftene i henholdsvis $V$ og $W$.

\subsection{Kvadratiske former og positive matricer}
En kvadratisk form kan skrives
\[
Q(x)=x^{\mathsf T}Kx.
\]
Kun den symmetriske del af $K$ bidrager, da $x^{\mathsf T}(K-K^{\mathsf T})x=0$. For en reel symmetrisk matrix $K$ gælder: $Q(x)>0$ for alle $x\ne0$ præcis når alle egenværdier af $K$ er positive. Denne karakterisering viser, hvorfor egenværdier bruges til at afgøre, om et stationært punkt er et stabilt lokalt minimum i en lineær eller kvadratisk model.

\section{Fysisk sammenhæng: koblede små svingninger}
\label{sec:normalmoder}
For små forskydninger $q$ omkring en ligevægt får man ofte en Lagrangefunktion på formen
\[
L=\frac12\dot q^{\mathsf T}M\dot q-\frac12q^{\mathsf T}Kq,
\]
hvor $M$ er en symmetrisk positiv definit massematrix, og $K$ er en symmetrisk stivhedsmatrix (Hessianen af potentialet i ligevægten). Euler--Lagrange giver
\[
M\ddot q+Kq=0.
\]
Ved ansatsen $q(t)=a\cos(\omega t+\varphi)$ fremkommer
\[
(K-\omega^2M)a=0,\qquad
\boxed{\det(K-\omega^2M)=0.}
\]
Dette er et \emph{generaliseret egenværdiproblem}. Det reduceres kun direkte til $Ka=\omega^2a$, når $M=I$ i de valgte enheder og koordinater. Hvis $K$ er positiv semidefinit, er de fundne $\omega^2$ ikke-negative; en nulmode kan svare til en fri bevægelsesretning. En negativ egenværdi af det relevante massenormaliserede problem giver i stedet en eksponentielt voksende mode.

Når $M$ er positiv definit, kan man indføre de massenormaliserede koordinater med $u=M^{1/2}a$ og den symmetriske matrix
\[
B=M^{-1/2}KM^{-1/2},
\qquad Bu=\omega^2u.
\]
Så kan den almindelige spektralsætning anvendes på $B$. Dette er forbindelsen mellem lineær algebra, andenordens differentialligninger og stabilitet i analytisk mekanik.

\subsection{Hvorfor superposition virker}
Da bevægelsesligningen er lineær, er summen af løsninger også en løsning. Når massenormaliserede egenvektorer danner en basis, kan de almindelige initialbetingelser opdeles på naturlige moder; hver mode svinger med sin egen frekvens. Superposition gælder i \emph{den lineære model}; store udsving kan kræve højere ordens led i potentialet, hvor normalmoderne ikke længere udvikler sig uafhængigt.

\section{Forbindelser til senere bind}
De matematiske resultater her er forudsætninger for en række senere emner. Jacobimatricen i flervariabel analyse er den lokale lineære afbildning, der bedst approksimerer et differentiabelt koordinatskift. Fouriermetoder skriver funktioner i en basis af oscillationer; Sturm--Liouville-teori generaliserer egenværdiproblemer til differentialoperatorer. Tensorregning skelner mellem objekter og deres koordinatkomponenter og generaliserer sporet ved kontraktion.

Den røde tråd er derfor ikke, at nye emner er det samme som lineær algebra, men at de ofte \emph{bevarer dens grundstruktur} under mere generelle betingelser.

\backmatter
\formelregister
\end{document}
